\section{Source corrections, integration decisions, and further research}
\label{audit:section}

\subsection{Precise changes to the source status}

No newly false inherited mathematical theorem was found in the
proof dependencies used here. The proposed changes are extensions
of valid theorems and corrections to research-status lists.
They should be applied with their original attribution intact.

\begin{enumerate}
\item In the ordered-boundary report's Section~9, replace the
statement that complete gap-six coverage is open by a reference
to Theorem~\ref{gs:main}, explicitly recording its
computation-assisted proof. Retain the broader all-pairs question.
The source's stronger sign conjecture
$f_b>0$ for every $b\ge7$ is not proved by a nonzero finite-field
residue. We establish positivity for $b\ge2048$ and nonvanishing
at all earlier admissible indices; we make no additional finite
sign claim.
\item In the bilinear determinant report's Section~10, mark the
complex rational-weight degree--rank question solved by
Theorem~\ref{cr:degree}. Broaden the generic rational-weight
degree statement to complex coefficients. Preserve the real
eventual-sign hypotheses in the source's sharper numerical
exception bound; that bound does not follow from the new proof.
\item In the lattice-contact report's Question~Q9 and the final
paragraph of its independent-shift section, mark the fixed
positive unequal-rate collision question solved by
Theorems~\ref{urc:constant-theorem} and \ref{urc:full-theorem}.
Keep nonlinear paths and ratios tending to zero or infinity open.
\item Do not reopen the already proved smaller-order-ten and
smaller-order-twelve cases. Do not relabel the arbitrary mixed
even-tail product calculus or the existing Stieltjes spectral-face
results as new consequences of this article.
\end{enumerate}

Several implementation conventions deserve explicit retention.
In the late Bernoulli coefficient matrix, a lower binomial index
of zero is omitted, rather than evaluated as $\binom n0=1$.
Changing that convention changes the theorem's matrix. In the
complex rank proof the relevant bilinear operation is ordinary
transposition; replacing it by conjugate transposition would
change the rational algebraic problem. In the Stieltjes section
the constant depends on the stated small parameter; critical
Laurent quotients must be removed jointly before taking spectral
derivatives. These are guards against incorrect integrations of
the results, not allegations that the cited source theorems made
those mistakes.

The manuscript's old rejected $S_8$ candidate should remain
rejected, and its revised $S_8$ and short $S_6$ candidates should
retain their conjectural status. Nothing in the present rank,
Hahn, or collision arguments implies those period equalities.
A small numerical residual does not supply their missing proof.

\subsection{A proposed research program}

The next questions below are selected for their potential to
produce exact identities or structural theorems. Each has a
concrete starting point supplied by the article.

\begin{enumerate}[label=\textbf{Q\arabic*.},leftmargin=*]
\item \textbf{Complete the gap-eight family.}
For half-gap $m=4$, make the inherited four-mode Gamma basis
effective. Two modes have a vanishing leading imaginary phase.
Derive a uniform two-term estimate for each, control the
four-by-four limiting determinant, and obtain an explicit finite
completion threshold. The target is nonvanishing for every
$b\ge5$, not simply a larger numerical table.

\item \textbf{Replace the gap-six finite completion by a structural proof.}
Find a recurrence invariant or a variation-diminishing theorem
which proves $f_b\ne0$ at every admissible $b$. A fixed-sign
assertion from the first index cannot work because the initial
sign transition is real. The stronger source conjecture
$f_b>0$ for $b\ge7$ would also settle the sign issue left
unresolved by the modular certificate.

\item \textbf{Connect fixed-gap and fixed-smaller-order regimes.}
The source has eventual theorems in both directions, and the
present proof makes one gap ray effective. Develop a
two-parameter determinant expansion uniform when both indices
increase. The limiting reciprocal-square nodes become less
separated as the gap grows; merely reusing fixed-gap constants
does not cover the intervening region.

\item \textbf{Evaluate the Gamma normalization determinants.}
The initial determinant is nonzero by an integral alternant
argument. Seek explicit Gamma, hypergeometric, or recurrence
evaluations for its normalization $K_m$ at $m=2,3,\ldots$.
Such evaluations would turn the recurrence certificates into
more compact identities rather than just nonvanishing tests.

\item \textbf{Classify rationally summable polynomial subspaces.}
For a given rational weight $R$, classify the spaces $V$ on
which every discrete-residue bilinear form of $Rfg$ vanishes.
Theorem~\ref{cr:primitive} proves nondegeneracy of every rational
primitive once $V$ is known; it does not classify these spaces.
This asks for a concrete bridge between rational summability
and symmetric pencils with varying kernels.

\item \textbf{Find intrinsic complex cutoff-exception invariants.}
Express the degree or factorization of the exceptional-cutoff
polynomial directly from the pole orders, residues, and integer
dispersions of $R$, without constructing the full determinant.
The coefficient-gcd criterion is exact, but a residue-level
description could explain and predict the later nonreal example
in Section~\ref{cr:section}.

\item \textbf{Extend the adjacent-pole compiler to several poles.}
For $R(n)=\sum_{\ell=1}^d c_\ell/(n+a_\ell)$, determine when
the highest harmonic coefficient reduces to a determinant of
fixed size as $r$ varies. For integer-separated poles the
compressed measure is a polynomial transform, suggesting a
Hahn or multiple-orthogonal representation. The difficult
part is to classify endpoint pole cancellations and exceptional
residue varieties, rather than invoke the general Christoffel
formula alone.

\item \textbf{Classify when classical polylogarithms suffice.}
Theorem~\ref{hp:all-r-generator} gives an entire arbitrary-rank
family using only one dilogarithm. Determine necessary and
sufficient conditions on a rational weight for its cutoff
generators to remain in a fixed classical-polylogarithm space.
Generic harmonic degree supplies an upper-level clue but is
not by itself a minimal-depth theorem.

\item \textbf{Evaluate the Euler primitives of the ensemble family.}
For the explicit $\mathcal Z_N$, determine all endpoint values
$\sum_{M\ge1}\mathcal Z_{M-1}/M^s$ in their convergence range,
and their alternating companions. Equation
\eqref{hp:euler-reduction} generalizes to a small list of
linear and quadratic Euler sums. The question is when these
reduce to ordinary zeta products and when irreducible higher
weight components first survive.

\item \textbf{Resolve genuinely different collision scales.}
Derive the full expansion for
$a=\lambda\epsilon^\alpha$ and $b=\mu\epsilon^\beta$,
$\alpha\ne\beta$. The fixed-ratio theorem cannot simply be
used with an $\epsilon$-dependent ratio without new uniform
analysis. Its beta kernel and finite contact polynomial
identify the terms which must enter a transition formula.

\item \textbf{Develop three-factor collision constants.}
For three independently shifted Stieltjes derivatives,
products of two singular factors with one analytic coefficient
may affect the constant. Determine the finite local contact
tensor and the appropriate multivariable beta kernel before
claiming a reduction to the two-factor alphabet.

\item \textbf{Identify invariant ordered combinations and complex jumps.}
Classify linear combinations of derivative products whose
collision constants depend only on the rate ratio, or are
independent of both rates. Then continue the rates in a
common complex sector and compute logarithmic jumps across
the boundary cuts. At fixed orders the first task is an
explicit finite polynomial problem in the contact coefficients.
\end{enumerate}

The strongest unresolved global targets remain the unrestricted
moment-prefix conjecture and the manuscript's distinguished
Gaussian period identities. The present article supplies three
completed local research programs and concrete next steps toward
broader statements; it does not infer those global targets from
finite evidence.

